\documentclass[10pt,a4paper]{amsart}
\usepackage[margin=19mm]{geometry}
\usepackage{amsmath,amssymb,mathtools}
\usepackage[scr=rsfs,cal=euler]{mathalfa}
\usepackage{xfrac}
\usepackage[unicode,hidelinks,bookmarks=false]{hyperref}
\newcommand{\ie}{i.e.\ }
\newcommand{\cf}{cf.\ }
\newcommand{\resp}{respectively\ }
\newcommand{\Subsection}{\subsection}
\providecommand{\nobreakdash}{\nobreak\hbox{-}}
\setlength{\emergencystretch}{2em}
\multlinegap0pt
\usepackage{bm}
\usepackage{bbm}
\usepackage{dsfont}
\usepackage{xspace}
\usepackage{cancel}
\usepackage{mathtools}
\usepackage{amsthm}
\newtheorem{theorem}{Theorem}
\newtheorem{proposition}[theorem]{Proposition}
\newcommand{\F}{\mathbf F}
\newcommand{\GL}{\operatorname{GL}}
\newcommand{\GSp}{\operatorname{GSp}}
\newcommand{\PGSp}{\operatorname{PGSp}}
\newcommand{\PP}{\mathbf P}
\newcommand{\ind}{\operatorname{ind}}
\title[Counting number fields with symplectic Galois group]{Counting number fields with symplectic Galois group}
\author{Anwesh Ray}
\date{Source: arXiv:2609.23093v1 (CC BY 4.0)}
\begin{document}
\maketitle

\noindent\emph{Source and licence.} Counting number fields with symplectic Galois group. Version arXiv:2609.23093v1 (CC BY 4.0), distributed under the Creative Commons Attribution 4.0 International licence, as recorded in the arXiv licence metadata for this exact version and shown on the abstract page https://arxiv.org/abs/2609.23093v1. Full rights notice: licenses/arxiv-2609.23093-CC-BY-4.0.txt.\par\smallskip

\noindent\textit{Editorial scope.} Counted unit: the Proposition whose source label is \texttt{prop:malle-invariants-natural} in the article (source0.tex lines 906--922, source label \texttt{prop:malle-invariants-natural}), delivered together with the article's own complete original proof of it (source0.tex lines 924--1010, one \texttt{proof} environment of 87 lines). No mathematics is written, repaired or re-derived: statement and proof are the article's own text line for line. Required context retained uncounted: none. Every definition, display and label of those lines is retained unchanged. Omitted from this package and not redistributed: 1--905, 923--923, 1011--1041. Nothing inside a retained excerpt is omitted or altered: every retained body is delivered exactly as the article's lines are written, so no omission receipt of any kind accompanies this package. Citations: the retained excerpts carry no citation command at all, so no bibliography entry of the article is transcribed and no reference is redistributed. Reference adaptation: none was needed and none was made; every reference inside the delivered unit resolves inside it, so no unresolved reference or citation is produced. Printed slips kept literal: no printed word, symbol, hypothesis or number of the counted statement or of its proof was altered. Layout and class: the article's own class is \texttt{amsart} with packages including geometry, amsmath, amssymb, amsthm, mathtools, enumitem, xcolor, tikz-cd, hyperref, orcidlink, mathrsfs, microtype, biblatex; the wrapper is the lane's pinned \texttt{amsart} editorial wrapper at 10pt on A4, which restates the theorem-like environments the retained text needs and adds the article's own macro definitions that the retained text needs (\texttt{\textbackslash F}, \texttt{\textbackslash GL}, \texttt{\textbackslash GSp}, \texttt{\textbackslash PGSp}, \texttt{\textbackslash PP}, \texttt{\textbackslash ind}), with extra wrapper packages (bm, bbm, dsfont, xspace, cancel, mathtools). The delivered numbering is the unit's own. Assets: no retained span contains a figure environment, a graphics-inclusion command, a PDF-page inclusion, a table or a TikZ picture; no image file is copied into this package, the package declares no asset dependency and no image file is redistributed with it. Licence: version arXiv:2609.23093v1 of this article is distributed under the Creative Commons Attribution 4.0 International licence, as recorded in the arXiv licence metadata for this exact version and shown on the abstract page https://arxiv.org/abs/2609.23093v1. A full rights notice is delivered at licenses/arxiv-2609.23093-CC-BY-4.0.txt. Characters: every retained excerpt is pure ASCII, so the delivered engine, XeLaTeX with its default Latin Modern text font, reports no missing character.
\begin{proposition}\label{prop:malle-invariants-natural}
Let $\pi_{\mathrm{vec}}$ be the natural action of
$\GSp_{2n}(\F_\ell)$ on $V\setminus\{0\}$, and let
$\pi_{\mathrm{proj}}$ be the natural action of
$\PGSp_{2n}(\F_\ell)$ on $\PP(V)$.  If $n\geq2$, then
\[
 a_{\pi_{\mathrm{vec}}}\bigl(\GSp_{2n}(\F_\ell)\bigr)
 =\frac{\ell^{2n-2}(\ell^2-1)}{2}
\]
and
\[
 a_{\pi_{\mathrm{proj}}}\bigl(\PGSp_{2n}(\F_\ell)\bigr)
 =\frac{(\ell+1)(\ell^{2n-2}-1)}{2}.
\]
For $n=1$, the corresponding values are
$\ell(\ell-1)/2$ and $(\ell-1)/2$, respectively.
\end{proposition}

\begin{proof}
We first record a reduction which will be used in both actions.  If a
nonidentity permutation $g$ has composite order, choose a nontrivial
power $h\in\langle g\rangle$ of prime order.  Every
$\langle g\rangle$-orbit is a union of $\langle h\rangle$-orbits, so
$\ind(h)\leq\ind(g)$.  Thus the minimum defining the Malle invariant is
attained by an element of prime order.  If such an element has prime
order $r$ on a set $\Omega$, then every orbit has size $1$ or $r$, and hence $\ind(g)=\frac{r-1}{r}(|\Omega|-|\operatorname{Fix}_\Omega(g)|)$.

Consider first the vector action and assume $n\geq2$.  Put $q=\ell$.
If $g$ has order $q$, then $g$ is unipotent.  Since $g\neq1$, its fixed
subspace has dimension at most $2n-1$, so $\ind(g)\geq q^{2n-2}(q-1)^2$.
For $q\geq3$ this is at least
$q^{2n-2}(q^2-1)/2$, with equality only when $q=3$ and the fixed
subspace is a hyperplane.

Now suppose that $g$ has prime order $r\neq q$, and write $\mu(g)$ for
its similitude multiplier.  The element is semisimple.  If
$\mu(g)=1$, then the $1$-eigenspace is a nondegenerate symplectic
subspace. The symplectic pairing pairs the $\lambda$- and
$\lambda^{-1}$-eigenspaces, so the $1$-eigenspace is orthogonal to all
other eigenspaces and the restriction of the pairing to it is
nondegenerate.  Since $g$ is nontrivial, this fixed space has dimension at most $2n-2$. Hence $\ind(g)\geq \frac12(q^{2n}-q^{2n-2})=\frac{q^{2n-2}(q^2-1)}2$.
If $\mu(g)\neq1$, then the fixed space is totally isotropic, because
for fixed vectors $v,w$ one has
$\langle v,w\rangle=\mu(g)\langle v,w\rangle$.  Its dimension is therefore at most $n$, and since $n\geq2$ the same lower bound follows from $\ind(g)\geq\frac12(q^{2n}-q^n)\geq\frac12(q^{2n}-q^{2n-2})$.
The bound is attained by a symplectic involution which is $-1$ on a
nondegenerate symplectic plane and $1$ on its symplectic orthogonal
complement.  This proves the vector formula for $n\geq2$.

For $n=1$, a nontrivial semisimple element of multiplier one has no
nonzero fixed vector, whereas an element with nontrivial multiplier can
have a one-dimensional fixed space.  The involution
$\operatorname{diag}(1,-1)\in\GL_2(\F_q)=\GSp_2(\F_q)$ realizes this
maximum and has index $q(q-1)/2$.  Unipotent elements have index
$(q-1)^2$, which is larger for odd $q$.  Hence the stated $n=1$ vector
value follows.

We next consider the projective action.  Again it is enough to consider a
projective element $\overline g$ of prime order $r$.  If it fixes a
projective point, choose a lift $g\in\GSp(V)$ and scale it so that one
fixed vector has eigenvalue $1$.  Since $\overline g^r=1$, this scaling
makes $g^r=1$; thus $g$ itself may be taken to have order $r$.  If
$\overline g$ has no fixed point, the resulting permutation index is
larger than the bounds below, so there is nothing to prove.

If $r=q$, then $g$ is unipotent and has only the eigenvalue $1$.  Its
fixed vector space has dimension at most $2n-1$, so the number of fixed projective points is at most $(q^{2n-1}-1)/(q-1)$. It follows that $\ind(\overline g)\geq q^{2n-2}(q-1)$.
For $n\geq2$ this is at least
$(q+1)(q^{2n-2}-1)/2$.

Suppose next that $r\neq q$.  Then $g$ is semisimple.  A projective
point fixed by $\overline g$ is an $\F_q$-rational eigenline of $g$.
Let $E_\lambda$ denote the eigenspace for an eigenvalue
$\lambda\in\F_q^\times$.  The similitude relation shows that
$E_\lambda$ pairs nondegenerately with $E_{\mu/\lambda}$, where
$\mu$ is the multiplier of $g$.  Consequently, if
$\lambda^2\neq\mu$, the two eigenspaces have the same dimension, while
if $\lambda^2=\mu$, the restriction of the alternating form to
$E_\lambda$ is nondegenerate and hence $\dim E_\lambda$ is even.  In
particular a non-scalar similitude has no eigenspace of dimension
$2n-1$.  Among all collections of rational eigenspaces with total
dimension at most $2n$ and with no eigenspace of dimension $2n-1$, the
number of eigenlines is therefore at most $(q^{2n-2}-1)/(q-1)+(q^2-1)/(q-1)$.
Indeed, the function $d\mapsto(q^d-1)/(q-1)$ is strictly convex in
$d$, so the number of lines is maximized by concentrating the available
dimension into blocks of dimensions $2n-2$ and $2$.  Hence
\[
 \ind(\overline g)
 \geq
 \frac12\left(
 \frac{q^{2n}-1}{q-1}
 -\frac{q^{2n-2}-1}{q-1}
 -\frac{q^2-1}{q-1}\right)
 =\frac{(q+1)(q^{2n-2}-1)}2.
\]
Equality is attained by the projective image of the same symplectic
involution which is $-1$ on a symplectic plane and $1$ on its
orthogonal complement.  This proves the projective formula for
$n\geq2$.

Finally, when $n=1$, a nontrivial projective semisimple element fixes at
most two points of $\PP^1(\F_q)$, and a split involution fixes exactly
two.  Its index is therefore $(q-1)/2$.  A projective unipotent element
fixes one point and has index $q-1$, so the involution gives the
minimum.  This completes the proof.
\end{proof}

\end{document}
